\section{Polynomial implementation and derandomization}
\label{sec:foundations-derandomization}

The analysis has so far considered a uniformly random common quantile.
We now show how to find a schedule at least as good as the expected
schedule by checking only polynomially many priorities. The LP has
polynomially many variables and constraints and rational coefficients.
An optimal rational solution of polynomial encoding length can therefore
be found in polynomial time \cite{Schrijver1986}. The algorithm uses the original-job
distributions only. Virtual jobs are needed for the proof, not for this
computation.

\begin{lemma}[Finitely many priority orders]
\label{lem:finite-priority-orders}
For a rational LP solution of polynomial encoding length in the input size,
one can construct in polynomial time a
partition of $(0,1)$ into $O(n^3)$ open intervals and their endpoints
such that the original-job priority order is constant in every open
interval. A rational point of polynomial encoding length can be chosen
in each interval.
\end{lemma}
\begin{proof}
For each job $i$ and each slot $s$ with $x_{i,s}>0$, the inverse
distribution is affine on $(F_i(s-1),F_i(s))$:
\[
 M_i(\alpha)=s-1+\frac{\alpha-F_i(s-1)}{x_{i,s}}.
\]
The endpoints $F_i(s)$ partition the quantile domain into at most
$n$ such pieces for each job. For a fixed pair $i,k$, merge the two
sets of endpoints. This gives $O(n)$ intervals, on each of which
$M_i-M_k$ is affine. If this difference has one isolated zero in an
interval, add it to the partition. If it is identically zero, the fixed
tie order determines the comparison throughout that interval. Otherwise
its sign is constant on each interval left after adding its zero, if any.
The union of the endpoints and all isolated zeros
therefore fixes every pairwise comparison on each resulting open
interval.

There are $O(n^2)$ pairs and $O(n)$ pieces per pair, for
$O(n^3)$ breakpoints. Each breakpoint is
obtained by rational arithmetic from LP entries of polynomial encoding
length. A midpoint between adjacent distinct breakpoints is rational and
also has polynomial encoding length. Constructing these points and
sorting them takes polynomial time.
\end{proof}

\begin{lemma}[Deterministic selection]
\label{lem:deterministic-selection}
The algorithm can find in polynomial time an original-job schedule with
cost at most the expected cost of the random-quantile schedule.
\end{lemma}
\begin{proof}
For every interval from Lemma~\ref{lem:finite-priority-orders}, run the
list rule on its chosen representative, then keep the cheapest schedule.
The same priority order gives the same schedule at every point of the
interval. The breakpoints themselves form a finite set and hence have
probability zero. Consequently the expected cost is the weighted average
of the finitely many representative costs, weighted by interval length.
Their minimum is no larger than that average. Each schedule is
constructed in polynomial time, and there are $O(n^3)$ candidates.
\end{proof}

\section{The approximation guarantee}

The following theorem combines the per-job bound with deterministic
selection.

\begin{theorem}[Approximation guarantee]
\label{thm:final}
There is a deterministic polynomial-time
\[
 \left(1+\sqrt2-0.0288468\right)
 \text{-approximation}
\]
for $P\mid\mathrm{prec},p_j=1\mid\sum_{j\in\mathcal O} w_j C_j$,
with arbitrary precedence directed acyclic graphs, nonnegative rational weights, and an
input-dependent number of identical machines. The ratio is approximately
$2.3853667624$.
\end{theorem}
\begin{proof}
If $m\ge n$, the depth schedule returned by Algorithm~\ref{alg:rounding}
is optimal by Section~\ref{sec:foundations-lp}. Assume $1\le m<n$.
Apply Lemma~\ref{thm:meta-symbolic} to an optimal LP solution,
multiply by the nonnegative weights, and sum over all original jobs. The
LP optimum is a lower bound on the scheduling optimum. The approximation
factor is positive, and Lemma~\ref{lem:deterministic-selection} gives a
schedule whose cost is no greater than the expected cost.
\end{proof}

\section{Conclusion}

We have presented a common-quantile rounding analysis for precedence
constrained unit-time scheduling and a parameter choice proving
a strict improvement over the $1+\sqrt2$ bound. The argument connects
priority-band progress to a scalar excess inequality and then couples
otherwise separate losses through a signed moment. The algorithm uses
the original jobs and a polynomial number of priority orders. Virtual capacity padding serves only the analysis.

The certificate establishes a feasible parameter choice rather than an
optimal choice within the analysis. Improving the scalar bounds or combining
several certificates may yield a stronger guarantee.

\appendix
\section{Parameter certificate}
\label{sec:meta-parameters}

\subsection{A family of sufficient certificates}

Consider real parameters $U,\ell,\kappa,Q,\lambda,\nu$ satisfying the band
conditions \eqref{eq:meta-band-domain}, the majorant
\eqref{eq:meta-majorant}, the defect condition
\eqref{eq:meta-defect-domain}, and the low-crossing conditions
\eqref{eq:meta-low-parameters}. Assume $0<U<1$, $Q>0$, and
$Q\cdot\kappa\ge c$.
Let $0<\delta\le c+1$ satisfy
\begin{equation}
\begin{aligned}
 \delta\cdot\max\bigl\{&1,\ Q-\kappa\cdot(1-\ell)-\nu+\lambda\cdot c,\\
 &2\cdot(2-U)\cdot(1-\ell)/c-\kappa\cdot(1-\ell)-\nu-\lambda\cdot(c+4),\\
 &\lambda\cdot(c+1)-\kappa\cdot(1-\ell)-\nu\bigr\}\\
 &\le\lambda-(c+2)\cdot[\kappa\cdot(1-\ell)+\nu]\\
 &\quad+\min\{0,(1-\ell)\cdot(1-U)+\kappa\cdot(1-\ell)+\nu-\lambda\cdot(c+1)\}.
\end{aligned}
\label{eq:general-certificate}
\end{equation}

\begin{proposition}[Parametric guarantee]
For $1\le m<n$, every original target $j$, and every feasible LP solution,
\[
 \mathbb E[\widetilde C_j]\le(c+2-\delta)\cdot C_j^{\mathrm{LP}}.
\]
Consequently, there is a deterministic polynomial-time
$(c+2-\delta)$-approximation for the weighted problem.
\end{proposition}
\begin{proof}
The scalar charges and coefficient collection in the proof of
Lemma~\ref{thm:meta-symbolic} use only the assumptions above. The maximum
in \eqref{eq:general-certificate} is at least one and bounds the three
coefficient deficits used in that proof. The corresponding factors are
nonnegative and sum to at most the gap by \eqref{eq:meta-base-gap}. The same absorption therefore
shows that this maximum, multiplied by
$(c+2)\cdot C_j^{\mathrm{LP}}-\mathbb E[\widetilde C_j]$, is at least
\[
 \bigl[(1-\ell)\cdot(1-U)+\kappa\cdot(1-\ell)+\nu-\lambda\cdot(c+1)\bigr]\cdot a_j
 +\bigl[\lambda-(c+2)\cdot(\kappa\cdot(1-\ell)+\nu)\bigr]\cdot C_j^{\mathrm{LP}}.
\]
By Lemma~\ref{lem:lp-support} and the assignment constraint,
\[
 C_j^{\mathrm{LP}}=\sum_{s=a_j}^{n}s\cdot x_{j,s}\ge a_j.
\]
If the first
coefficient is negative, replacing $a_j$ by $C_j^{\mathrm{LP}}$ only
decreases the expression. If it is nonnegative, dropping the first term
does the same. In either case, \eqref{eq:general-certificate} completes
the per-job bound. Weighting and deterministic selection give the stated
approximation when $m<n$, as in the proof of Theorem~\ref{thm:final}.
If $m\ge n$, the depth schedule is optimal and $c+2-\delta\ge1$.
\end{proof}

One can seek the supremum of $\delta$ subject to these conditions. Each
feasible choice gives a sufficient certificate. This family fixes the upper
quantile endpoint at $1$ and uses the same coefficient $Q$ in both scalar cases. We certify one
feasible choice below and make no claim that it is optimal.

\subsection{A feasible choice}

We now verify the parameter inequalities and coefficient bounds used in
Lemma~\ref{thm:meta-symbolic}. Choose
\begin{center}
\begin{tabular}{c|r@{\qquad}c|r}
 Parameter & Value & Parameter & Value\\ \hline
 $U$ & $0.411323$ & $\ell$ & $0.491938$\\
 $\kappa$ & $0.108423$ & $Q$ & $3.820348$\\
 $\lambda$ & $0.252884$ & $\nu$ & $0.003462$
\end{tabular}
\end{center}
These values give
\[
 \ell-\kappa=0.383515,\qquad 1-\ell=0.508062.
\]
In particular, $0<\ell-\kappa<U<\ell<1$, $0<\kappa<1/2$, and
$1-\ell>\kappa$. Moreover, with $c=\sqrt2-1$,
\[
 Q\cdot\kappa-c>2.8\cdot10^{-8}.
\]

\subsection{Checking the majorant}

Using $(c+2)\cdot c=1$ and $(c+2)\cdot c^2=c$, define the quadratic
\[
 R(u)=-\frac{Q\cdot(c+2)}2\cdot u^2+(Q-\lambda)\cdot u
      +c\cdot(\lambda-Q/2)-\nu.
\]
The residual of \eqref{eq:meta-majorant} is
\[
 R(u)+(1-\ell)\cdot(u-U)_++\Gamma(u).
\]
On $[0,\ell-\kappa]$ both additional terms vanish. The derivative of $R$
is decreasing and satisfies
\[
 R'(\ell-\kappa)>0.03025,\qquad R(\ell-\kappa)<-0.0000447.
\]
Thus the residual is negative on this interval.

For the remaining intervals, the following table gives the coefficients
of the polynomial added to $R(u)$ for each branch of $\Gamma$. In each row,
the rightmost entry is a strict upper bound on that polynomial plus $R(u)$
for all real $u$. The first two rows use $u\in[\ell-\kappa,U]$, the next
two use $u\in[U,\ell]$, and the last uses $u\in[\ell,1]$.
\begin{center}
\small
\begin{tabular}{c|c|c|c|r}
 Branch & coefficient of $u^2$ & coefficient of $u$ & constant & upper bound\\ \hline
 Middle quadratic & $-1/2$ & $\ell-\kappa$ & $-(\ell-\kappa)^2/2$ & $-10^{-9}$\\
 Middle linear & $0$ & $1/4$ & $-\ell/4$ & $-0.0228$\\
 Middle quadratic & $-1/2$ & $1-\kappa$ & $-(\ell-\kappa)^2/2-(1-\ell)\cdot U$ & $-5.82\cdot10^{-8}$\\
 Middle linear & $0$ & $5/4-\ell$ & $-\ell/4-(1-\ell)\cdot U$ & $-0.00758$\\
 Upper linear & $0$ & $2-\kappa-\ell$ & $-(1-\kappa)\cdot\ell-(1-\ell)\cdot U$ & $-6.89\cdot10^{-8}$
\end{tabular}
\end{center}
To check a row, write its sum with $R(u)$ as $-a\cdot u^2+b\cdot u+d$,
where $a>0$. Completing the square shows that its maximum over the real
line is $d+b^2/(4\cdot a)$. For example, the first row has
\[
 a=\frac{Q\cdot(c+2)+1}{2},\quad
 b=Q-\lambda+\ell-\kappa,\quad
 d=c\cdot(\lambda-Q/2)-\nu-\frac{(\ell-\kappa)^2}{2}.
\]
Substituting the stated parameters yields
\[
 4\cdot a\cdot\left(d+\frac{b^2}{4\cdot a}+10^{-9}\right)
 <5.51959502-3.90294321\cdot\sqrt2<0.
\]
The last inequality follows by squaring the positive terms, since
$5.51959502^2<2\cdot3.90294321^2$. Therefore the first row's maximum is
less than $-10^{-9}$. The other rows use the same calculation. Both middle
branches are bounded separately, so their maximum is bounded as well.
This proves \eqref{eq:meta-majorant}.

\subsection{Remaining inequalities}

Direct substitution gives
\[
 \frac{Q\cdot(1-c)}2+\lambda\cdot c-\nu-[1+\kappa\cdot(2-\ell)]
 >\frac1{20}.
\]
This proves \eqref{eq:meta-defect-domain}. The two inequalities in
\eqref{eq:meta-low-parameters} also hold with slack:
\[
 1+\lambda-\frac{1-\ell}{U}>\frac1{100},\qquad
 Q\cdot c-\bigl[1-\kappa\cdot(1-\ell)-\nu+\lambda\cdot c\bigr]>\frac12.
\]

Finally, the coefficient of $a_j$ in the proof of
Lemma~\ref{thm:meta-symbolic} is positive:
\[
 (1-\ell)\cdot(1-U)+\kappa\cdot(1-\ell)+\nu-\lambda\cdot(c+1)
 >3\cdot10^{-8}.
\]
Direct substitution also gives
\begin{equation}
\begin{aligned}
 \max\bigl\{&1,\ Q-\kappa\cdot(1-\ell)-\nu+\lambda\cdot c,\\
 &2\cdot(2-U)\cdot(1-\ell)/c-\kappa\cdot(1-\ell)-\nu-\lambda\cdot(c+4),\\
 &\lambda\cdot(c+1)-\kappa\cdot(1-\ell)-\nu\bigr\}
 <3.8665485
\end{aligned}
\end{equation}
and
\[
 \lambda-(c+2)\cdot[\kappa\cdot(1-\ell)+\nu]
 -0.0288468\cdot3.8665485>10^{-8}.
\]
These bounds complete the numerical step in Lemma~\ref{thm:meta-symbolic}.
They also verify \eqref{eq:general-certificate} for $\delta=0.0288468$.
